\documentclass[11pt, a4paper]{article}
% ── Shared preamble for all RTOS lab documents ────────────────────────────────
% Usage: % ── Shared preamble for all RTOS lab documents ────────────────────────────────
% Usage: % ── Shared preamble for all RTOS lab documents ────────────────────────────────
% Usage: \input{../preamble}  (from each labNN/ directory)
% ──────────────────────────────────────────────────────────────────────────────

% ── Packages ──────────────────────────────────────────────────────────────────
\usepackage[a4paper, top=2.5cm, bottom=2.5cm, left=2.5cm, right=2.5cm, headheight=28pt]{geometry}
\usepackage[utf8]{inputenc}
\usepackage[T1]{fontenc}
\usepackage{lmodern}
\usepackage{booktabs}
\usepackage{array}
\usepackage{tabularx}
\usepackage{longtable}
\usepackage{multirow}
\usepackage{colortbl}
\usepackage{xcolor}
\usepackage{titlesec}
\usepackage{enumitem}
\usepackage{hyperref}
\usepackage{fancyhdr}
\usepackage{parskip}
\usepackage{listings}
\usepackage[most]{tcolorbox}
\usepackage{fontawesome5}
\usepackage{microtype}
\usepackage{graphicx}
\usepackage{amsmath}

% ── Colors (KMUTNB brand) ──────────────────────────────────────────────────────
\definecolor{kmutnbBlue}{RGB}{0,56,116}
\definecolor{kmutnbGold}{RGB}{186,146,36}
\definecolor{lightGray}{RGB}{245,245,245}
\definecolor{midGray}{RGB}{200,200,200}
\definecolor{codeGray}{RGB}{248,248,248}
\definecolor{codeFrame}{RGB}{220,220,220}
\definecolor{tipteal}{RGB}{0,105,92}
\definecolor{warnAmber}{RGB}{121,85,0}
\definecolor{checkGreen}{RGB}{27,94,32}

% ── Section formatting ─────────────────────────────────────────────────────────
\titleformat{\section}
  {\large\bfseries\color{kmutnbBlue}}
  {\thesection.}{0.5em}{}
  [\color{kmutnbGold}\titlerule]

\titleformat{\subsection}
  {\normalsize\bfseries\color{kmutnbBlue}}
  {\thesubsection.}{0.5em}{}

\titleformat{\subsubsection}
  {\small\bfseries\color{kmutnbBlue}}
  {\thesubsubsection.}{0.5em}{}

\titlespacing*{\section}{0pt}{1.4ex plus .2ex}{0.6ex plus .1ex}
\titlespacing*{\subsection}{0pt}{1.0ex plus .2ex}{0.4ex plus .1ex}

% ── Listings (C and bash) ──────────────────────────────────────────────────────
\lstdefinestyle{cstyle}{
  language        = C,
  basicstyle      = \ttfamily\small,
  keywordstyle    = \bfseries\color{kmutnbBlue},
  commentstyle    = \itshape\color{gray},
  stringstyle     = \color{tipteal},
  numberstyle     = \tiny\color{midGray},
  numbers         = left,
  numbersep       = 6pt,
  stepnumber      = 1,
  backgroundcolor = \color{codeGray},
  frame           = single,
  framerule       = 0.4pt,
  rulecolor       = \color{codeFrame},
  breaklines      = true,
  showstringspaces= false,
  tabsize         = 4,
  xleftmargin     = 1.5em,
  xrightmargin    = 0.5em,
  aboveskip       = 0.8em,
  belowskip       = 0.8em,
}

\lstdefinestyle{bashstyle}{
  language        = bash,
  basicstyle      = \ttfamily\small,
  keywordstyle    = \color{kmutnbBlue},
  commentstyle    = \itshape\color{gray},
  backgroundcolor = \color{black!5},
  frame           = single,
  framerule       = 0.4pt,
  rulecolor       = \color{codeFrame},
  breaklines      = true,
  showstringspaces= false,
  xleftmargin     = 1.5em,
  xrightmargin    = 0.5em,
  aboveskip       = 0.6em,
  belowskip       = 0.6em,
}

\lstdefinestyle{textstyle}{
  basicstyle      = \ttfamily\footnotesize,
  backgroundcolor = \color{black!4},
  frame           = single,
  framerule       = 0.4pt,
  rulecolor       = \color{codeFrame},
  breaklines      = true,
  xleftmargin     = 1.5em,
  xrightmargin    = 0.5em,
  aboveskip       = 0.6em,
  belowskip       = 0.6em,
}

% ── Callout boxes ──────────────────────────────────────────────────────────────
\tcbuselibrary{skins, breakable}

\newtcolorbox{tipbox}[1][]{
  enhanced, breakable,
  colback=tipteal!8, colframe=tipteal, coltitle=white,
  fonttitle=\bfseries\small, title={\faLightbulb\quad Tip#1},
  left=6pt, right=6pt, top=4pt, bottom=4pt,
  boxrule=0.8pt, arc=3pt,
}

\newtcolorbox{warnbox}[1][]{
  enhanced, breakable,
  colback=warnAmber!8, colframe=kmutnbGold, coltitle=white,
  fonttitle=\bfseries\small, title={\faExclamationTriangle\quad Note#1},
  left=6pt, right=6pt, top=4pt, bottom=4pt,
  boxrule=0.8pt, arc=3pt,
}

\newtcolorbox{checkpointbox}[1][]{
  enhanced,
  colback=kmutnbBlue!6, colframe=kmutnbBlue, coltitle=white,
  fonttitle=\bfseries\small, title={\faCheckCircle\quad Checkpoint#1},
  left=6pt, right=6pt, top=4pt, bottom=4pt,
  boxrule=0.8pt, arc=3pt,
}

\newtcolorbox{questionbox}[1][]{
  enhanced, breakable,
  colback=kmutnbGold!10, colframe=kmutnbGold, coltitle=white,
  fonttitle=\bfseries\small, title={\faQuestion\quad Question#1},
  left=6pt, right=6pt, top=4pt, bottom=4pt,
  boxrule=0.8pt, arc=3pt,
}

% ── Checkbox list ──────────────────────────────────────────────────────────────
\newlist{checklist}{itemize}{1}
\setlist[checklist]{label=\faSquare[regular], leftmargin=1.5em, noitemsep}

% ── Misc helpers ───────────────────────────────────────────────────────────────
\newcommand{\file}[1]{\texttt{\small #1}}
\newcommand{\api}[1]{\texttt{\small #1}}

% ── Title block macro ─────────────────────────────────────────────────────────
% Usage: \labtitle{03}{Aperiodic Servers \& Producer--Consumer}{2}{CLO 1 \& CLO 2}
\newcommand{\labtitle}[4]{%
  \begin{center}
    {\color{kmutnbBlue}\rule{\linewidth}{2pt}}\\[6pt]
    {\LARGE\bfseries\color{kmutnbBlue} Laboratory #1}\\[4pt]
    {\large\bfseries #2}\\[2pt]
    {\normalsize Real-Time Operating Systems \enspace\textbullet\enspace
     M.Eng.\ ECE \enspace\textbullet\enspace KMUTNB}\\[8pt]
    {\color{kmutnbGold}\rule{\linewidth}{1pt}}
  \end{center}
  \vspace{0.3cm}
  \begin{tabular}{@{}p{3.8cm} p{10cm}@{}}
    \textbf{Platform}       & NXP FRDM-MCXN236 (ARM Cortex-M33 @ 150\,MHz) \\[2pt]
    \textbf{RTOS}           & FreeRTOS v10.6 \\[2pt]
    \textbf{Toolchain}      & ARM GNU Toolchain (\texttt{arm-none-eabi-gcc}) + Make \\[2pt]
    \textbf{Estimated time} & #3 hours \\[2pt]
    \textbf{CLO mapping}    & #4 \\
  \end{tabular}
  \vspace{0.4cm}
}
  (from each labNN/ directory)
% ──────────────────────────────────────────────────────────────────────────────

% ── Packages ──────────────────────────────────────────────────────────────────
\usepackage[a4paper, top=2.5cm, bottom=2.5cm, left=2.5cm, right=2.5cm, headheight=28pt]{geometry}
\usepackage[utf8]{inputenc}
\usepackage[T1]{fontenc}
\usepackage{lmodern}
\usepackage{booktabs}
\usepackage{array}
\usepackage{tabularx}
\usepackage{longtable}
\usepackage{multirow}
\usepackage{colortbl}
\usepackage{xcolor}
\usepackage{titlesec}
\usepackage{enumitem}
\usepackage{hyperref}
\usepackage{fancyhdr}
\usepackage{parskip}
\usepackage{listings}
\usepackage[most]{tcolorbox}
\usepackage{fontawesome5}
\usepackage{microtype}
\usepackage{graphicx}
\usepackage{amsmath}

% ── Colors (KMUTNB brand) ──────────────────────────────────────────────────────
\definecolor{kmutnbBlue}{RGB}{0,56,116}
\definecolor{kmutnbGold}{RGB}{186,146,36}
\definecolor{lightGray}{RGB}{245,245,245}
\definecolor{midGray}{RGB}{200,200,200}
\definecolor{codeGray}{RGB}{248,248,248}
\definecolor{codeFrame}{RGB}{220,220,220}
\definecolor{tipteal}{RGB}{0,105,92}
\definecolor{warnAmber}{RGB}{121,85,0}
\definecolor{checkGreen}{RGB}{27,94,32}

% ── Section formatting ─────────────────────────────────────────────────────────
\titleformat{\section}
  {\large\bfseries\color{kmutnbBlue}}
  {\thesection.}{0.5em}{}
  [\color{kmutnbGold}\titlerule]

\titleformat{\subsection}
  {\normalsize\bfseries\color{kmutnbBlue}}
  {\thesubsection.}{0.5em}{}

\titleformat{\subsubsection}
  {\small\bfseries\color{kmutnbBlue}}
  {\thesubsubsection.}{0.5em}{}

\titlespacing*{\section}{0pt}{1.4ex plus .2ex}{0.6ex plus .1ex}
\titlespacing*{\subsection}{0pt}{1.0ex plus .2ex}{0.4ex plus .1ex}

% ── Listings (C and bash) ──────────────────────────────────────────────────────
\lstdefinestyle{cstyle}{
  language        = C,
  basicstyle      = \ttfamily\small,
  keywordstyle    = \bfseries\color{kmutnbBlue},
  commentstyle    = \itshape\color{gray},
  stringstyle     = \color{tipteal},
  numberstyle     = \tiny\color{midGray},
  numbers         = left,
  numbersep       = 6pt,
  stepnumber      = 1,
  backgroundcolor = \color{codeGray},
  frame           = single,
  framerule       = 0.4pt,
  rulecolor       = \color{codeFrame},
  breaklines      = true,
  showstringspaces= false,
  tabsize         = 4,
  xleftmargin     = 1.5em,
  xrightmargin    = 0.5em,
  aboveskip       = 0.8em,
  belowskip       = 0.8em,
}

\lstdefinestyle{bashstyle}{
  language        = bash,
  basicstyle      = \ttfamily\small,
  keywordstyle    = \color{kmutnbBlue},
  commentstyle    = \itshape\color{gray},
  backgroundcolor = \color{black!5},
  frame           = single,
  framerule       = 0.4pt,
  rulecolor       = \color{codeFrame},
  breaklines      = true,
  showstringspaces= false,
  xleftmargin     = 1.5em,
  xrightmargin    = 0.5em,
  aboveskip       = 0.6em,
  belowskip       = 0.6em,
}

\lstdefinestyle{textstyle}{
  basicstyle      = \ttfamily\footnotesize,
  backgroundcolor = \color{black!4},
  frame           = single,
  framerule       = 0.4pt,
  rulecolor       = \color{codeFrame},
  breaklines      = true,
  xleftmargin     = 1.5em,
  xrightmargin    = 0.5em,
  aboveskip       = 0.6em,
  belowskip       = 0.6em,
}

% ── Callout boxes ──────────────────────────────────────────────────────────────
\tcbuselibrary{skins, breakable}

\newtcolorbox{tipbox}[1][]{
  enhanced, breakable,
  colback=tipteal!8, colframe=tipteal, coltitle=white,
  fonttitle=\bfseries\small, title={\faLightbulb\quad Tip#1},
  left=6pt, right=6pt, top=4pt, bottom=4pt,
  boxrule=0.8pt, arc=3pt,
}

\newtcolorbox{warnbox}[1][]{
  enhanced, breakable,
  colback=warnAmber!8, colframe=kmutnbGold, coltitle=white,
  fonttitle=\bfseries\small, title={\faExclamationTriangle\quad Note#1},
  left=6pt, right=6pt, top=4pt, bottom=4pt,
  boxrule=0.8pt, arc=3pt,
}

\newtcolorbox{checkpointbox}[1][]{
  enhanced,
  colback=kmutnbBlue!6, colframe=kmutnbBlue, coltitle=white,
  fonttitle=\bfseries\small, title={\faCheckCircle\quad Checkpoint#1},
  left=6pt, right=6pt, top=4pt, bottom=4pt,
  boxrule=0.8pt, arc=3pt,
}

\newtcolorbox{questionbox}[1][]{
  enhanced, breakable,
  colback=kmutnbGold!10, colframe=kmutnbGold, coltitle=white,
  fonttitle=\bfseries\small, title={\faQuestion\quad Question#1},
  left=6pt, right=6pt, top=4pt, bottom=4pt,
  boxrule=0.8pt, arc=3pt,
}

% ── Checkbox list ──────────────────────────────────────────────────────────────
\newlist{checklist}{itemize}{1}
\setlist[checklist]{label=\faSquare[regular], leftmargin=1.5em, noitemsep}

% ── Misc helpers ───────────────────────────────────────────────────────────────
\newcommand{\file}[1]{\texttt{\small #1}}
\newcommand{\api}[1]{\texttt{\small #1}}

% ── Title block macro ─────────────────────────────────────────────────────────
% Usage: \labtitle{03}{Aperiodic Servers \& Producer--Consumer}{2}{CLO 1 \& CLO 2}
\newcommand{\labtitle}[4]{%
  \begin{center}
    {\color{kmutnbBlue}\rule{\linewidth}{2pt}}\\[6pt]
    {\LARGE\bfseries\color{kmutnbBlue} Laboratory #1}\\[4pt]
    {\large\bfseries #2}\\[2pt]
    {\normalsize Real-Time Operating Systems \enspace\textbullet\enspace
     M.Eng.\ ECE \enspace\textbullet\enspace KMUTNB}\\[8pt]
    {\color{kmutnbGold}\rule{\linewidth}{1pt}}
  \end{center}
  \vspace{0.3cm}
  \begin{tabular}{@{}p{3.8cm} p{10cm}@{}}
    \textbf{Platform}       & NXP FRDM-MCXN236 (ARM Cortex-M33 @ 150\,MHz) \\[2pt]
    \textbf{RTOS}           & FreeRTOS v10.6 \\[2pt]
    \textbf{Toolchain}      & ARM GNU Toolchain (\texttt{arm-none-eabi-gcc}) + Make \\[2pt]
    \textbf{Estimated time} & #3 hours \\[2pt]
    \textbf{CLO mapping}    & #4 \\
  \end{tabular}
  \vspace{0.4cm}
}
  (from each labNN/ directory)
% ──────────────────────────────────────────────────────────────────────────────

% ── Packages ──────────────────────────────────────────────────────────────────
\usepackage[a4paper, top=2.5cm, bottom=2.5cm, left=2.5cm, right=2.5cm, headheight=28pt]{geometry}
\usepackage[utf8]{inputenc}
\usepackage[T1]{fontenc}
\usepackage{lmodern}
\usepackage{booktabs}
\usepackage{array}
\usepackage{tabularx}
\usepackage{longtable}
\usepackage{multirow}
\usepackage{colortbl}
\usepackage{xcolor}
\usepackage{titlesec}
\usepackage{enumitem}
\usepackage{hyperref}
\usepackage{fancyhdr}
\usepackage{parskip}
\usepackage{listings}
\usepackage[most]{tcolorbox}
\usepackage{fontawesome5}
\usepackage{microtype}
\usepackage{graphicx}
\usepackage{amsmath}

% ── Colors (KMUTNB brand) ──────────────────────────────────────────────────────
\definecolor{kmutnbBlue}{RGB}{0,56,116}
\definecolor{kmutnbGold}{RGB}{186,146,36}
\definecolor{lightGray}{RGB}{245,245,245}
\definecolor{midGray}{RGB}{200,200,200}
\definecolor{codeGray}{RGB}{248,248,248}
\definecolor{codeFrame}{RGB}{220,220,220}
\definecolor{tipteal}{RGB}{0,105,92}
\definecolor{warnAmber}{RGB}{121,85,0}
\definecolor{checkGreen}{RGB}{27,94,32}

% ── Section formatting ─────────────────────────────────────────────────────────
\titleformat{\section}
  {\large\bfseries\color{kmutnbBlue}}
  {\thesection.}{0.5em}{}
  [\color{kmutnbGold}\titlerule]

\titleformat{\subsection}
  {\normalsize\bfseries\color{kmutnbBlue}}
  {\thesubsection.}{0.5em}{}

\titleformat{\subsubsection}
  {\small\bfseries\color{kmutnbBlue}}
  {\thesubsubsection.}{0.5em}{}

\titlespacing*{\section}{0pt}{1.4ex plus .2ex}{0.6ex plus .1ex}
\titlespacing*{\subsection}{0pt}{1.0ex plus .2ex}{0.4ex plus .1ex}

% ── Listings (C and bash) ──────────────────────────────────────────────────────
\lstdefinestyle{cstyle}{
  language        = C,
  basicstyle      = \ttfamily\small,
  keywordstyle    = \bfseries\color{kmutnbBlue},
  commentstyle    = \itshape\color{gray},
  stringstyle     = \color{tipteal},
  numberstyle     = \tiny\color{midGray},
  numbers         = left,
  numbersep       = 6pt,
  stepnumber      = 1,
  backgroundcolor = \color{codeGray},
  frame           = single,
  framerule       = 0.4pt,
  rulecolor       = \color{codeFrame},
  breaklines      = true,
  showstringspaces= false,
  tabsize         = 4,
  xleftmargin     = 1.5em,
  xrightmargin    = 0.5em,
  aboveskip       = 0.8em,
  belowskip       = 0.8em,
}

\lstdefinestyle{bashstyle}{
  language        = bash,
  basicstyle      = \ttfamily\small,
  keywordstyle    = \color{kmutnbBlue},
  commentstyle    = \itshape\color{gray},
  backgroundcolor = \color{black!5},
  frame           = single,
  framerule       = 0.4pt,
  rulecolor       = \color{codeFrame},
  breaklines      = true,
  showstringspaces= false,
  xleftmargin     = 1.5em,
  xrightmargin    = 0.5em,
  aboveskip       = 0.6em,
  belowskip       = 0.6em,
}

\lstdefinestyle{textstyle}{
  basicstyle      = \ttfamily\footnotesize,
  backgroundcolor = \color{black!4},
  frame           = single,
  framerule       = 0.4pt,
  rulecolor       = \color{codeFrame},
  breaklines      = true,
  xleftmargin     = 1.5em,
  xrightmargin    = 0.5em,
  aboveskip       = 0.6em,
  belowskip       = 0.6em,
}

% ── Callout boxes ──────────────────────────────────────────────────────────────
\tcbuselibrary{skins, breakable}

\newtcolorbox{tipbox}[1][]{
  enhanced, breakable,
  colback=tipteal!8, colframe=tipteal, coltitle=white,
  fonttitle=\bfseries\small, title={\faLightbulb\quad Tip#1},
  left=6pt, right=6pt, top=4pt, bottom=4pt,
  boxrule=0.8pt, arc=3pt,
}

\newtcolorbox{warnbox}[1][]{
  enhanced, breakable,
  colback=warnAmber!8, colframe=kmutnbGold, coltitle=white,
  fonttitle=\bfseries\small, title={\faExclamationTriangle\quad Note#1},
  left=6pt, right=6pt, top=4pt, bottom=4pt,
  boxrule=0.8pt, arc=3pt,
}

\newtcolorbox{checkpointbox}[1][]{
  enhanced,
  colback=kmutnbBlue!6, colframe=kmutnbBlue, coltitle=white,
  fonttitle=\bfseries\small, title={\faCheckCircle\quad Checkpoint#1},
  left=6pt, right=6pt, top=4pt, bottom=4pt,
  boxrule=0.8pt, arc=3pt,
}

\newtcolorbox{questionbox}[1][]{
  enhanced, breakable,
  colback=kmutnbGold!10, colframe=kmutnbGold, coltitle=white,
  fonttitle=\bfseries\small, title={\faQuestion\quad Question#1},
  left=6pt, right=6pt, top=4pt, bottom=4pt,
  boxrule=0.8pt, arc=3pt,
}

% ── Checkbox list ──────────────────────────────────────────────────────────────
\newlist{checklist}{itemize}{1}
\setlist[checklist]{label=\faSquare[regular], leftmargin=1.5em, noitemsep}

% ── Misc helpers ───────────────────────────────────────────────────────────────
\newcommand{\file}[1]{\texttt{\small #1}}
\newcommand{\api}[1]{\texttt{\small #1}}

% ── Title block macro ─────────────────────────────────────────────────────────
% Usage: \labtitle{03}{Aperiodic Servers \& Producer--Consumer}{2}{CLO 1 \& CLO 2}
\newcommand{\labtitle}[4]{%
  \begin{center}
    {\color{kmutnbBlue}\rule{\linewidth}{2pt}}\\[6pt]
    {\LARGE\bfseries\color{kmutnbBlue} Laboratory #1}\\[4pt]
    {\large\bfseries #2}\\[2pt]
    {\normalsize Real-Time Operating Systems \enspace\textbullet\enspace
     M.Eng.\ ECE \enspace\textbullet\enspace KMUTNB}\\[8pt]
    {\color{kmutnbGold}\rule{\linewidth}{1pt}}
  \end{center}
  \vspace{0.3cm}
  \begin{tabular}{@{}p{3.8cm} p{10cm}@{}}
    \textbf{Platform}       & NXP FRDM-MCXN236 (ARM Cortex-M33 @ 150\,MHz) \\[2pt]
    \textbf{RTOS}           & FreeRTOS v10.6 \\[2pt]
    \textbf{Toolchain}      & ARM GNU Toolchain (\texttt{arm-none-eabi-gcc}) + Make \\[2pt]
    \textbf{Estimated time} & #3 hours \\[2pt]
    \textbf{CLO mapping}    & #4 \\
  \end{tabular}
  \vspace{0.4cm}
}


% ── Header / Footer ────────────────────────────────────────────────────────────
\pagestyle{fancy}
\fancyhf{}
\fancyhead[L]{\small\color{kmutnbBlue}\textbf{KMUTNB -- Faculty of Engineering}}
\fancyhead[R]{\small\color{kmutnbBlue}Dept.\ of Electrical and Computer Engineering}
\fancyfoot[L]{\small\color{kmutnbBlue}Real-Time Operating Systems -- M.Eng.\ ECE}
\fancyfoot[C]{\small\thepage}
\fancyfoot[R]{\small\color{kmutnbBlue}Lab 04}
\renewcommand{\headrulewidth}{0.4pt}
\renewcommand{\headrule}{\hbox to\headwidth{%
  \color{kmutnbGold}\leaders\hrule height\headrulewidth\hfill}}
\renewcommand{\footrulewidth}{0.4pt}
\renewcommand{\footrule}{\hbox to\headwidth{%
  \color{midGray}\leaders\hrule height\footrulewidth\hfill}}

\hypersetup{
  colorlinks = true,
  linkcolor  = kmutnbBlue,
  urlcolor   = kmutnbBlue,
  citecolor  = kmutnbBlue,
  pdftitle   = {Lab 04 -- Priority Inversion: See It, Fix It},
  pdfauthor  = {KMUTNB RTOS Course},
}

% ══════════════════════════════════════════════════════════════════════════════
\begin{document}

\labtitle{04}{Priority Inversion: See It, Fix It}{2--3}{CLO 3 \& CLO 4}

% ── Objectives ────────────────────────────────────────────────────────────────
\section{Objectives}

By completing this lab you will be able to:

\begin{enumerate}[noitemsep, leftmargin=*]
  \item \textbf{Reproduce} the classic three-task priority inversion scenario on real hardware.
  \item \textbf{Observe} the inversion window in a SEGGER SystemView trace and measure its
        duration.
  \item \textbf{Quantify} the difference between theoretical worst-case response time with and
        without PIP.
  \item \textbf{Fix} the inversion by replacing a binary semaphore with a FreeRTOS mutex.
  \item \textbf{Verify} the fix --- $\tau_H$'s response time must not exceed $\tau_L$'s
        critical-section length.
\end{enumerate}

% ── Prerequisites ─────────────────────────────────────────────────────────────
\section{Prerequisites}

\begin{checklist}
  \item Lab 03 complete --- SystemView setup working.
  \item SEGGER SystemView running and connected to the board.
  \item ARM GNU Toolchain and \texttt{pyocd} working.
\end{checklist}

% ── Background ────────────────────────────────────────────────────────────────
\section{Background}

Priority inversion occurs when a high-priority task ($\tau_H$) blocks on a resource held
by a low-priority task ($\tau_L$), and a medium-priority task ($\tau_M$) preempts
$\tau_L$ --- causing $\tau_H$ to wait for $\tau_M$, even though $\tau_M$ shares no
resource with $\tau_H$.

\textbf{Without PIP:}
Response time of $\tau_H \geq$ duration of $\tau_M$'s execution + $\tau_L$'s remaining
critical section. This can be arbitrarily long.

\textbf{With PIP (FreeRTOS mutex):}
When $\tau_H$ blocks, $\tau_L$ inherits $\tau_H$'s priority. $\tau_M$ cannot preempt
$\tau_L$. $\tau_H$'s blocking is bounded to at most one $\tau_L$ critical section.

% ── Project Structure ─────────────────────────────────────────────────────────
\section{Project Structure}

\begin{lstlisting}[style=textstyle]
lab04_priority_inversion/
  Makefile
  linker_MCXN236.ld
  src/
    main.c               <- task creation, kernel object init
    inversion_tasks.c/.h <- tauH, tauM, tauL implementations
    FreeRTOSConfig.h
    uart.h / uart.c
  FreeRTOS-Kernel/
\end{lstlisting}

\subsection{Task Parameters}

\renewcommand{\arraystretch}{1.3}
\begin{tabular}{@{} p{2.5cm} c c p{5cm} @{}}
  \toprule
  \textbf{Task} & \textbf{Priority} & \textbf{Period} & \textbf{Critical section} \\
  \midrule
  $\tau_H$ (high)   & 3 & 200\,ms & Takes shared resource for 5\,ms \\
  $\tau_M$ (medium) & 2 & 150\,ms & Runs 30\,ms of computation (no shared resource) \\
  $\tau_L$ (low)    & 1 & 300\,ms & Holds shared resource for 40\,ms \\
  \bottomrule
\end{tabular}
\renewcommand{\arraystretch}{1}

% ── Part A ─────────────────────────────────────────────────────────────────────
\section{Part A --- Reproduce the Inversion (Binary Semaphore)}

\subsection{Step 1 --- Implement the Three Tasks}

In \file{inversion\_tasks.c}, use a \textbf{binary semaphore}
(\texttt{xSemaphoreCreateBinary}, pre-given to \texttt{pdTRUE}) as the shared resource
guard:

\begin{lstlisting}[style=cstyle]
void vTaskH(void *pv)   /* prio 3 */
{
    TickType_t xLastWake = xTaskGetTickCount();
    for (;;) {
        vTaskDelayUntil(&xLastWake, pdMS_TO_TICKS(200));
        TickType_t t_arrive = xTaskGetTickCount();
        uart_printf("[H] arrived at %lu ms\r\n", t_arrive);

        xSemaphoreTake(xSharedResource, portMAX_DELAY);
        TickType_t t_acquire = xTaskGetTickCount();
        uart_printf("[H] acquired at %lu ms  (waited %lu ms)\r\n",
                    t_acquire, t_acquire - t_arrive);

        vTaskDelay(pdMS_TO_TICKS(5));   /* critical section */
        xSemaphoreGive(xSharedResource);
        uart_printf("[H] released resource\r\n");
    }
}

void vTaskM(void *pv)   /* prio 2 */
{
    TickType_t xLastWake = xTaskGetTickCount();
    for (;;) {
        vTaskDelayUntil(&xLastWake, pdMS_TO_TICKS(150));
        uart_printf("[M] running (no resource held)\r\n");
        TickType_t start = xTaskGetTickCount();
        while ((xTaskGetTickCount() - start) < pdMS_TO_TICKS(30)) {}
        uart_printf("[M] done\r\n");
    }
}

void vTaskL(void *pv)   /* prio 1 */
{
    TickType_t xLastWake = xTaskGetTickCount();
    for (;;) {
        vTaskDelayUntil(&xLastWake, pdMS_TO_TICKS(300));
        xSemaphoreTake(xSharedResource, portMAX_DELAY);
        uart_printf("[L] acquired resource -- holding for 40 ms\r\n");
        TickType_t start = xTaskGetTickCount();
        while ((xTaskGetTickCount() - start) < pdMS_TO_TICKS(40)) {}
        xSemaphoreGive(xSharedResource);
        uart_printf("[L] released resource\r\n");
    }
}
\end{lstlisting}

\begin{warnbox}
  Tasks are offset at startup so that $\tau_L$ acquires the resource first, $\tau_H$
  arrives while $\tau_L$ holds it, and $\tau_M$ preempts $\tau_L$ shortly after. This is
  handled in \file{main.c} using initial \texttt{vTaskDelay} offsets.
\end{warnbox}

\subsection{Step 2 --- Build and Flash}

\begin{lstlisting}[style=bashstyle]
make setup && make && make flash
\end{lstlisting}

\subsection{Step 3 --- Capture UART Output}

Watch for lines where $\tau_H$'s \texttt{waited X ms} value is much larger than the
40\,ms $\tau_L$ critical section alone.

Expected inversion output:
\begin{lstlisting}[style=textstyle]
[L] acquired resource -- holding for 40 ms
[H] arrived at 100 ms
[M] running (no resource held)
[M] done
[L] released resource
[H] acquired resource at 170 ms  (waited 70 ms)
\end{lstlisting}

\subsection{Step 4 --- Capture SystemView Trace}

Open SystemView, start recording, let the system run for at least 3 seconds, then stop.
In the trace, find an inversion event:
\begin{itemize}[noitemsep, leftmargin=*]
  \item Identify the moment $\tau_H$ blocks on the semaphore.
  \item Identify the moment $\tau_M$ preempts $\tau_L$.
  \item Measure the \textbf{inversion window} --- the time $\tau_H$ waits while $\tau_M$
        is running.
\end{itemize}

\begin{checkpointbox}[~A]
  Annotated SystemView screenshot with the inversion window marked and measured.
\end{checkpointbox}

% ── Part B ─────────────────────────────────────────────────────────────────────
\section{Part B --- Theoretical Analysis}

Before fixing the system, compute the worst-case response time analytically.

Without PIP, $\tau_H$'s response time can be bounded as:
\[
  R_H \leq C_H + B_H + \sum_{j \in \mathrm{hp}(H)} \left\lceil \frac{R_H}{T_j} \right\rceil C_j
\]
where $B_H$ is the \textbf{blocking term}. Without PIP, $B_H$ is unbounded ($\tau_M$ can
always preempt $\tau_L$).

With PIP, $B_H$ is bounded to the longest critical section of any lower-priority task
sharing the resource:
\[
  B_H = \max(\text{CS of lower-priority tasks}) = 40\,\mathrm{ms}
\]

\begin{questionbox}[~B1]
  Using the task parameters from the table, apply the RTA recurrence to compute
  $\tau_H$'s worst-case response time \textbf{with PIP}. Show your iteration steps.
\end{questionbox}

\begin{questionbox}[~B2]
  From your UART output in Part~A, what is the maximum \texttt{waited X ms} value you
  observed across multiple periods? Does it match the theoretical without-PIP bound?
\end{questionbox}

\begin{questionbox}[~B3]
  Compute total utilisation $U$ for the three-task set. Is it schedulable under RMS?
  Does priority inversion change the schedulability answer? (Hint: yes and no ---
  explain.)
\end{questionbox}

% ── Part C ─────────────────────────────────────────────────────────────────────
\section{Part C --- Fix with a Mutex}

\subsection{Step 1 --- Replace Binary Semaphore with Mutex}

In \file{main.c}, change one line:

\begin{lstlisting}[style=cstyle]
/* BEFORE (binary semaphore -- no PIP): */
xSharedResource = xSemaphoreCreateBinary();
xSemaphoreGive(xSharedResource);

/* AFTER (mutex -- PIP enabled): */
xSharedResource = xSemaphoreCreateMutex();
\end{lstlisting}

Also ensure \file{FreeRTOSConfig.h} has:
\begin{lstlisting}[style=cstyle]
#define configUSE_MUTEXES  1
\end{lstlisting}

No other code changes are needed --- \texttt{xSemaphoreTake} and \texttt{xSemaphoreGive}
have the same API for both mutexes and binary semaphores.

\subsection{Step 2 --- Rebuild and Flash}

\begin{lstlisting}[style=bashstyle]
make && make flash
\end{lstlisting}

\subsection{Step 3 --- Observe the Fix}

$\tau_H$'s \texttt{waited} time should now be $\leq 40$\,ms. $\tau_M$ cannot preempt
$\tau_L$ while $\tau_L$ holds the mutex.

Expected fixed output:
\begin{lstlisting}[style=textstyle]
[L] acquired resource -- holding for 40 ms
[H] arrived at 100 ms
[H] acquired resource at 140 ms  (waited 40 ms)
[M] running (no resource held)
[M] done
\end{lstlisting}

\subsection{Step 4 --- Capture the Fixed SystemView Trace}

In the fixed trace, $\tau_L$'s priority \textbf{spikes to $\tau_H$'s priority level}
when $\tau_H$ blocks (visible in SystemView as a priority change event).

\begin{checkpointbox}[~C]
  Annotated SystemView screenshot showing the priority inheritance event
  ($\tau_L$ running at $\tau_H$'s priority).
\end{checkpointbox}

% ── Part D ─────────────────────────────────────────────────────────────────────
\section{Part D --- Experiments}

\subsection{Experiment 1 --- Deep Nesting}

Add a \textbf{second} shared resource \texttt{xResource2} (also a mutex). Modify $\tau_L$
to hold both resources:

\begin{lstlisting}[style=cstyle]
xSemaphoreTake(xSharedResource, portMAX_DELAY);
vTaskDelay(pdMS_TO_TICKS(10));
xSemaphoreTake(xResource2, portMAX_DELAY);
/* ... critical section ... */
xSemaphoreGive(xResource2);
xSemaphoreGive(xSharedResource);
\end{lstlisting}

And $\tau_M$ also competes for \texttt{xResource2}.

\begin{questionbox}[~D1]
  FreeRTOS PIP documentation states that the current implementation propagates
  inheritance only \textbf{one level deep} --- it does not fully propagate chained
  inheritance. Describe a scenario with this two-resource design where $\tau_H$ still
  experiences unexpected delay even with mutexes.
\end{questionbox}

\subsection{Experiment 2 --- ISR Give}

Replace the \texttt{xSemaphoreGive} inside \texttt{vTaskL} with a give from a software
timer callback (simulating an ISR-side release).

\begin{questionbox}[~D2]
  \texttt{xSemaphoreGiveFromISR} does \textbf{not} trigger priority inheritance
  restoration. After $\tau_L$ ``gives'' from the ISR callback, does $\tau_L$'s priority
  return to its original level immediately? What does this tell you about using mutexes
  for resources released from ISR context?
\end{questionbox}

% ── Deliverables ───────────────────────────────────────────────────────────────
\section{Deliverables}

\renewcommand{\arraystretch}{1.3}
\begin{tabular}{@{} c p{4.5cm} p{7cm} @{}}
  \toprule
  \textbf{\#} & \textbf{Item} & \textbf{Details} \\
  \midrule
  1 & Part A UART capture   & At least one inversion event with \texttt{waited} annotated \\
  2 & Part A SystemView screenshot & Inversion window marked and measured \\
  3 & Part B analysis       & RTA calculation with and without PIP \\
  4 & Part C UART capture   & \texttt{waited} bounded to $\leq 40$\,ms \\
  5 & Part C SystemView screenshot & Priority inheritance event visible \\
  6 & Written answers       & Questions B1--B3, D1--D2 \\
  \bottomrule
\end{tabular}
\renewcommand{\arraystretch}{1}

% ── Key Concepts ───────────────────────────────────────────────────────────────
\section{Key Concepts Demonstrated}

\renewcommand{\arraystretch}{1.3}
\begin{tabular}{@{} p{7cm} p{6cm} @{}}
  \toprule
  \textbf{Concept} & \textbf{Where you saw it} \\
  \midrule
  Three-task priority inversion scenario      & Part A \\
  Inversion window measurement                & Part A SystemView trace \\
  RTA blocking term $B_H$                     & Part B analysis \\
  Mutex vs.\ binary semaphore for PIP         & Part C code change \\
  Priority inheritance in SystemView          & Part C trace \\
  Limits of FreeRTOS PIP (one-level)          & Experiment D1 \\
  \bottomrule
\end{tabular}
\renewcommand{\arraystretch}{1}

% ── Reference ─────────────────────────────────────────────────────────────────
\section{Reference}

\renewcommand{\arraystretch}{1.3}
\begin{tabular}{@{} p{5.5cm} p{7.5cm} @{}}
  \toprule
  \textbf{Resource} & \textbf{Relevant sections} \\
  \midrule
  Sha, Rajkumar \& Lehoczky (1990) &
    ``Priority Inheritance Protocols'' (\textit{IEEE Trans. Computers}) \\
  Buttazzo, \textit{Hard Real-Time Computing Systems} & Ch.~7.1--7.3 \\
  FreeRTOS docs & \texttt{xSemaphoreCreateMutex}, \texttt{configUSE\_MUTEXES} \\
  SEGGER SystemView manual & Priority display, event annotations \\
  \bottomrule
\end{tabular}
\renewcommand{\arraystretch}{1}

\end{document}
